% W4i47_Compiled.tex
% DFD 17-Agent Research Programme --- Iteration 47 (i47)
% ALL 17 AGENTS: Master Compiled Document
% Theme: Phase 0A Execution Crisis, Birefringence npi Escalation, Sigma mnu Resolution Paths, Scorecard 61/5/4 (5th iteration)
% Date: 2026-03-23

\documentclass[11pt,a4paper]{article}
\usepackage[utf8]{inputenc}
\usepackage{amsmath,amssymb,amsfonts,amsthm}
\usepackage{hyperref}
\usepackage{booktabs}
\usepackage{longtable}
\usepackage{geometry}
\usepackage{xcolor}
\usepackage{tcolorbox}
\geometry{margin=2.5cm}

\newtheorem{theorem}{Theorem}
\newtheorem{proposition}{Proposition}
\newtheorem{lemma}{Lemma}
\newtheorem{corollary}{Corollary}
\newtheorem{definition}{Definition}
\newtheorem{remark}{Remark}

\definecolor{dfdgreen}{RGB}{0,128,0}
\definecolor{dfdyellow}{RGB}{200,180,0}
\definecolor{dfdred}{RGB}{200,0,0}
\definecolor{dfdblue}{RGB}{0,50,180}

\newcommand{\GREEN}{\textcolor{dfdgreen}{\textbf{GREEN}}}
\newcommand{\YELLOW}{\textcolor{dfdyellow}{\textbf{YELLOW}}}
\newcommand{\RED}{\textcolor{dfdred}{\textbf{RED}}}
\newcommand{\NEW}{\textcolor{dfdblue}{\textbf{NEW}}}
\newcommand{\RESOLVED}{\textcolor{dfdgreen}{\textbf{RESOLVED}}}
\newcommand{\Tr}{\operatorname{Tr}}
\newcommand{\CP}{\mathbb{C}P}

\title{DFD i47: Compiled Master Document --- All 17 Agents\\[6pt]
\large Iteration 16/20 of Quantum Gravity Cycle (i32--i51)\\[6pt]
\normalsize Phase 0A Crisis $\cdot$ Birefringence $n\pi$ Shift $\cdot$ $\Sigma m_\nu$ Paths $\cdot$ Scorecard 61/5/4}
\author{DFD 17-Agent Research Programme}
\date{2026-03-23}

\begin{document}
\maketitle

\begin{abstract}
Iteration 47 is defined by three developments. First, the birefringence risk landscape has shifted: a new PRL paper identifies an $n\pi$ phase ambiguity meaning the true rotation angle could be much larger than $0.3^\circ$, while a separate PRL connects birefringence to the $\Sigma m_\nu$ tension, and a third PRL proposes an independent galaxy polarization probe. Second, the $\Sigma m_\nu$ tension has diversified: three independent resolution paths (spatial curvature, dynamical DE, birefringence) are now published, reducing the risk to DFD's 0.061~eV prediction. Third, the Phase 0A execution crisis deepens to six iterations without a single computation; Agent 14 correctly identifies the root cause (text-generation vs.\ numerical work) and recommends decoupling into a dedicated computational session. Additionally: the mixed anomaly is formally dropped after 7 iterations of deferral, ET appoints Michele Maggiore as spokesperson, the EMR feasibility study passes its continue-unless milestone, DESI dynamical DE evidence faces a robustness challenge, and wide binary tests continue converging toward GR. Scorecard: 61 GREEN, 5 YELLOW, 4 RED --- unchanged for the fifth consecutive iteration.
\end{abstract}

\tableofcontents
\newpage

%=============================================================================
\part{Executive Summary}
%=============================================================================

\section{i47 vs.\ i46 Comparison}

\begin{center}
\begin{tabular}{lll}
\toprule
\textbf{Aspect} & \textbf{i46} & \textbf{i47} \\
\midrule
Critical resolutions & 2 (b-quark, PRL partial) & 0 \\
Agenda changes & 0 & \textbf{1 (mixed anomaly DROPPED)} \\
New literature & $\sim$85 papers ($\sim$30 new) & $\sim$90+ papers ($\sim$35 new) \\
Scorecard & 61/5/4 & 61/5/4 (5th iteration frozen) \\
Phase 0A & Plan specified, unexecuted & \textbf{Root cause identified, decoupling recommended} \\
Birefringence & CS assessment performed & \textbf{$n\pi$ ambiguity + galaxy probe + $\Sigma m_\nu$ link} \\
$\Sigma m_\nu$ & FC tension emerging & \textbf{Three resolution paths identified} \\
Key external & JUNO first results, ET 3-site & \textbf{ET leadership, EMR feasibility, DESI debate} \\
Grade changes & 0 & \textbf{2 downgrades (Agent 3 A-$\to$B+, Agent 14 B$\to$C+)} \\
\bottomrule
\end{tabular}
\end{center}

\section{Top 5 New Developments}

\begin{enumerate}
\item \textbf{Birefringence $n\pi$ phase ambiguity} (PRL, Naokawa \& Namikawa, January 2026): The CMB EB correlation has an $n\pi$ phase degeneracy. The EB spectrum shape may resolve it, but the true angle could be $\gg 0.3^\circ$. Escalates birefringence risk for large-angle scenario. Simultaneously, a galaxy polarization probe (PRL 134, 2025) offers CMB-independent measurement with different systematics.

\item \textbf{$\Sigma m_\nu$ three resolution paths}: (a) Spatial curvature reduces tension from $2.59\sigma$ to $1.17\sigma$ (arXiv:2603.13208). (b) Dynamical DE relaxes bound to $< 0.163$ eV (arXiv:2603.10787). (c) Birefringence during reionization resolves negative mass preference (PRL, arXiv:2506.22999). DFD's 0.061 eV is safe in all three.

\item \textbf{Phase 0A root cause identified}: The iterative text-generation cycle cannot produce numerical computations. Phase 0A requires a dedicated computational session with Julia, SymBoltz.jl, and 12 uninterrupted hours. Six iterations of excellent planning, zero execution.

\item \textbf{Mixed anomaly formally dropped}: After 7 iterations of deferral (i40--i46), the mixed anomaly $c_1^{(Y)}$ computation is removed from the active agenda. Honest record: important open question, programme limitation.

\item \textbf{ET enters decisive phase}: Michele Maggiore appointed spokesperson. EMR feasibility passes continue-unless milestone. Bid book due December 2026. Site decision 2027.
\end{enumerate}


%=============================================================================
\part{Agent Grades: All 17}
%=============================================================================

\begin{center}
\begin{longtable}{clp{9cm}}
\toprule
\textbf{Agent} & \textbf{Grade} & \textbf{Key Result} \\
\midrule
\endfirsthead
\toprule
\textbf{Agent} & \textbf{Grade} & \textbf{Key Result} \\
\midrule
\endhead
1 ($\alpha$) & B+ & $n\pi$ birefringence ambiguity tracked. Birefringence--$\Sigma m_\nu$ link noted. \\
2 (Topology) & B+ & NCG spacetime algebra paper tracked. PRL impediment unchanged. \\
3 (Masses) & \textbf{B+} $\downarrow$ & Mass paper edit deferred 5 iters. $\Sigma m_\nu$ resolution paths analyzed. \\
4 (Spin$^c$) & B & $k_f$ gap remains. Phase 0B connection unchanged. \\
5 (Anomaly) & C+ & \textbf{Mixed anomaly DROPPED.} $\Sigma m_\nu$ three-path analysis. \\
6 (Dead patents) & C+ & Patent landscape static. \\
7 (ET/CE) & A- & Maggiore appointed. EMR continue-unless. IR1 confirmed. \\
8 (Laboratory) & A- & Th-229 programme consolidated. No new data. \\
9 (Particle) & B+ & Moriond 2026 fully digested. $HH \to bb\gamma\gamma$ limits. \\
10 (Astro) & B+ & Wide binary triple modelling supports GR. Gaia DR4 December 2026. \\
11 (CMB) & B+ & $n\pi$ birefringence. Galaxy polarization probe. Birefringence--$\Sigma m_\nu$ link. \\
12 (LSS) & B+ & DESI robustness debate tracked. Qualitative Euclid pre-reg recommended. \\
13 (Dark) & B & $\Sigma m_\nu$ resolution landscape. Scalar NSI risk to JUNO noted. \\
14 (Numerics) & \textbf{C+} $\downarrow$ & Phase 0A unexecuted for 6 iters. Root cause identified. gCAMB noted. \\
15 (Scorecard) & B & 61/5/4 frozen for 5 iterations. \\
16 (Cross-check) & A- & All claims verified. Process failures escalated. Circularity in birefringence--$\Sigma m_\nu$ noted. \\
17 (Synthesis) & A- & Honest trajectory. Phase 0A decoupling recommended. \\
\bottomrule
\end{longtable}
\end{center}

\textbf{Mean grade}: B+ (unchanged from i46).

\textbf{Grade changes} (i46 $\to$ i47):
\begin{itemize}
\item Agent 3: A- $\to$ \textbf{B+} (mass paper edit deferred 5 iterations).
\item Agent 14: B $\to$ \textbf{C+} (Phase 0A unexecuted 6 iterations, despite excellent plan quality).
\end{itemize}


%=============================================================================
\part{Scorecard: 70 Predictions}
%=============================================================================

\section{Summary}

\begin{center}
\begin{tabular}{lccc}
\toprule
& \textbf{\GREEN} & \textbf{\YELLOW} & \textbf{\RED} \\
\midrule
i46 & 61 & 5 & 4 \\
i47 & 61 & 5 & 4 \\
\textbf{Change} & \textbf{0} & \textbf{0} & \textbf{0} \\
\bottomrule
\end{tabular}
\end{center}

Fifth consecutive iteration at 61/5/4. The only internally controllable promotion ($A_L$) requires Phase 0A execution.

\section{RED Predictions (4)}

\begin{enumerate}
\item GW echo amplitude ($10^{-41}$): untestable
\item BMV enhancement ($2200\times$ QED): awaiting data
\item Lepton universality: anomalies dissolved, pending Run 3
\item $|\lambda-1|$ direct: no experiment sensitive
\end{enumerate}

\section{YELLOW Predictions (5)}

\begin{enumerate}
\item $\dot{\alpha}/\alpha$ at $10^{-19}$/yr: Th-229 by $\sim$2030 (on track)
\item EFE in galaxies: converging to GR, Gaia DR4 December 2026
\item $A_L$ lensing: $1.025 \pm 0.017$ vs.\ DFD $1.00$--$1.05$ (Phase 0A needed)
\item $\Sigma m_\nu$: $< 0.064$ eV vs.\ DFD $0.061$ eV (\textbf{risk reduced: three resolution paths})
\item Dynamical dark energy: DESI $2.8$--$4.2\sigma$ (\textbf{robustness debated})
\end{enumerate}


%=============================================================================
\part{Phase 0: Critical Path Update}
%=============================================================================

\section{Phase 0A: Root Cause Analysis and Decoupling Recommendation}

\begin{tcolorbox}[colback=red!5!white, colframe=red!60!black,
  title={Phase 0A: SIX ITERATIONS WITHOUT EXECUTION --- Root Cause and Fix}]

\textbf{Root cause}: The 17-agent iterative cycle generates LaTeX text. Phase 0A requires Julia code execution, package installation, and numerical computation. These are fundamentally different activities.

\textbf{Proposed fix}: Decouple Phase 0A into a \textbf{dedicated computational session}:
\begin{enumerate}
\item Schedule a standalone session (not an iteration cycle).
\item Install Julia 1.10+ and SymBoltz.jl.
\item Execute the 4-step, 12-hour plan specified in i46.
\item Report results back into the iterative cycle.
\end{enumerate}

\textbf{Deadline constraint}: Euclid pre-registration 20 June 2026 (88 days). If Phase 0A is not executed by i49, the programme will fail on its single most important deliverable.

\textbf{Status}: Plan A-level. Execution F-level.
\end{tcolorbox}

\section{Phase 0B: Framework Outlined, Unstarted}

Five parameters ($m_\psi$, $c_s^2$, $\pi_\psi$, photon/baryon coupling, initial conditions) remain underived. The $\psi$ perturbation theory derivation has not been attempted. Key literature identified (K-essence Boltzmann dynamics, $F(T)$ gravity codes, SymBoltz.jl custom submodels).


%=============================================================================
\part{Birefringence: Updated Risk Assessment}
%=============================================================================

\section{New Developments (i47)}

Three new PRL papers shift the birefringence landscape:

\begin{enumerate}
\item \textbf{$n\pi$ phase ambiguity} (PRL, January 2026): The true rotation angle has a degeneracy $\beta + n\pi$. The EB spectrum shape may constrain $n$, but results suggest the angle could be larger than $0.3^\circ$.

\item \textbf{Birefringence resolves negative $\Sigma m_\nu$} (PRL, October 2025, arXiv:2506.22999): Birefringence during reionization can simultaneously account for the EB signal, the elevated $\tau$, and the negative effective neutrino mass.

\item \textbf{Galaxy polarization probe} (PRL 134, 161001, 2025): Independent birefringence measurement via spiral galaxy radio polarization. SKA surveys could match CMB-S4 sensitivity with different systematics.
\end{enumerate}

\section{Updated Risk Matrix}

\begin{center}
\begin{tabular}{lp{3cm}p{3cm}p{3cm}}
\toprule
\textbf{Scenario} & \textbf{Birefringence} & \textbf{$\Sigma m_\nu$} & \textbf{DFD Impact} \\
\midrule
$\beta \approx 0.3^\circ$, confirmed & MEDIUM risk & Safe (relaxed bound) & Requires CS extension (Option 3) \\
$\beta \gg 0.3^\circ$ ($n > 0$) & \textbf{HIGH risk} & Very safe (large relaxation) & Only Option 1 survives easily \\
$\beta$ unconfirmed (systematic) & LOW risk & Tension remains & No action needed \\
\bottomrule
\end{tabular}
\end{center}

\textbf{Key insight from Agent 16}: The birefringence--$\Sigma m_\nu$ connection creates a coupled problem. Resolving $\Sigma m_\nu$ via birefringence requires DFD to accommodate birefringence. This is circular but potentially productive: a single CS extension could resolve both.

\textbf{Timeline}: Simons Observatory ($\sim$2027) and LiteBIRD ($\sim$2028) will be decisive.


%=============================================================================
\part{Falsification Risk Dashboard}
%=============================================================================

\begin{center}
\begin{tabular}{lcccl}
\toprule
\textbf{Observable} & \textbf{DFD} & \textbf{Data} & \textbf{Risk} & \textbf{i47 Change} \\
\midrule
$\Sigma m_\nu$ & 0.061 eV & $< 0.064$ eV & \YELLOW & \textbf{Risk reduced: 3 resolution paths} \\
$A_L$ & 1.00--1.05 & $1.025 \pm 0.017$ & \GREEN & No change \\
$R_{31}$ & TBD & --- & \YELLOW & No change \\
Birefringence & None & $\geq 0.3^\circ$ & \RED & \textbf{$n\pi$ ambiguity: HIGH if $n > 0$} \\
GW lensing & Never lensed & None & \GREEN & No change \\
$D_L^{\text{EM}}/D_L^{\text{GW}}$ & $\sim 1.31$ & N/A & \GREEN & No change \\
Mass ordering & Normal & $2.2\sigma$ NO & \GREEN & Accumulating \\
\bottomrule
\end{tabular}
\end{center}

\textbf{Birefringence remains the highest-risk item.} The $n\pi$ ambiguity increases uncertainty on the severity. $\Sigma m_\nu$ risk has decreased.


%=============================================================================
\part{New Literature: Complete List ($\sim$90+ papers across 17 agents)}
%=============================================================================

\section{Theory (Agents 1--5)}

\begin{enumerate}
\item PRL (2026) --- $n\pi$ phase ambiguity of cosmic birefringence. \NEW.
\item PRL (2025) --- Resolving negative $\Sigma m_\nu$ with birefringence (arXiv:2506.22999). \NEW.
\item PRL 134, 161001 (2025) --- Galaxy polarization birefringence probe. \NEW.
\item arXiv:2601.07350 (January 2026) --- Novel NCG spacetime algebra. \NEW.
\item arXiv:2603.10787 (March 2026) --- $\Sigma m_\nu$ with ACT DR6 + DESI DR2. \NEW.
\item arXiv:2603.13208 (March 2026) --- Negative masses and spatial curvature. \NEW.
\item arXiv:2602.05564 (February 2026) --- Scalar NSI risk to JUNO. \NEW.
\item arXiv:2512.08752 (December 2025) --- $\Sigma m_\nu$ in 4-parameter DE. \NEW.
\item Nature Communications (2025) --- $K_\alpha$ for Th-229. (continued)
\item MNRAS 528, 6550 (2024) --- Quasar $\alpha$ convergence. (continued)
\item JHEP (February 2025) --- Spectral torsion for Connes type operator. (continued)
\item JHEP (February 2026) --- LEFT two-loop RGEs. (continued)
\item arXiv:2511.14593 --- JUNO first results. (continued)
\item arXiv:2601.09791 --- Lessons from JUNO. (continued)
\end{enumerate}

\section{Observables (Agents 8--14)}

\begin{enumerate}
\item ATLAS $HH \to bb\gamma\gamma$ (March 2026) --- Record Higgs self-coupling limits. \NEW.
\item ATLAS $\tau \to 3\mu$ (March 2026) --- LFV search. \NEW.
\item arXiv:2504.07569 (April 2025) --- Wide binaries GR vs MOND with triple modelling. \NEW.
\item arXiv:2504.15222 (April 2025) --- Did DESI DR2 truly reveal dynamical DE? \NEW.
\item arXiv:2511.22512 (November 2025) --- Robust DDE evidence from DESI DR2. \NEW.
\item Nature Astronomy (2025) --- Dynamical DE in light of DESI DR2. \NEW.
\item arXiv:2601.04312 (January 2026) --- Neutrino decays explanation. \NEW.
\item arXiv:2603.03284 (March 2026) --- Decaying DM background vs perturbations. \NEW.
\item gCAMB (A\&C, 2026) --- GPU-accelerated Boltzmann solver. \NEW.
\item Nature (January 2026) --- Th-229 frequency reproducibility. (continued)
\item Euclid Prep.\ LXXXV (A\&A, March 2026) --- DR1 HO-WL. (continued)
\item SymBoltz.jl (A\&A, March 2026) --- Published. (continued)
\item Qu et al.\ (PRL 136, 2026) --- $A_L = 1.025 \pm 0.017$. (continued)
\end{enumerate}

\section{Cosmology and GW (Agents 6, 7, 15, 16, 17)}

\begin{enumerate}
\item ET Collaboration (March 2026) --- Maggiore appointed spokesperson. \NEW.
\item EMR feasibility (March 2026) --- Continue-unless milestone. \NEW.
\item GWTC-4.0 (arXiv:2508.18080) --- 218 events. (continued)
\item ET-SUnLab tender (February 2026) --- (continued)
\item Sardinia--Saxony cooperation (January 2026) --- (continued)
\item arXiv:2508.20721 --- O4a isotropic GW background. (continued)
\item LVK IR1 plans --- September--October 2026. (continued)
\end{enumerate}


%=============================================================================
\part{Action Items for i48}
%=============================================================================

\begin{center}
\begin{tabular}{clcc}
\toprule
\textbf{Priority} & \textbf{Action} & \textbf{Owner} & \textbf{Blocks} \\
\midrule
\textbf{CRITICAL} & \textbf{Schedule dedicated Phase 0A session} & Gary / Agent 14 & Everything \\
\textbf{HIGH} & Draft qualitative Euclid pre-registration & Agent 12 & Euclid deadline \\
\textbf{HIGH} & Execute mass paper edit or DROP & Agent 3 & Publication \\
MEDIUM & Track birefringence $n\pi$ developments & Agent 11 & Falsification risk \\
MEDIUM & Monitor $\Sigma m_\nu$ resolution papers & Agent 5/13 & Prediction viability \\
MEDIUM & Track DESI robustness debate & Agent 12 & DE interpretation \\
MEDIUM & Begin Phase 0B derivation ($\psi$ perturbation eqns) & Agent 14 & Full predictions \\
LOW & $k_f$ quantitative computation & Agent 2/4 & Mass exponents \\
\bottomrule
\end{tabular}
\end{center}

\textbf{Deadline reminder}: Euclid pre-registration 20 June 2026 (\textbf{88 days}). LHC Run 3 ends June 2026. ET-SUnLab tender 16 April 2026 (24 days). LVK IR1 September--October 2026. ET bid book December 2026. Gaia DR4 December 2026.

\textbf{Programme status}: Iteration 16/20 (80\% complete). Four iterations remain. Phase 0A execution in a dedicated session is the single most important action. The programme's success or failure hinges on this one deliverable.

\end{document}
